\documentclass[11pt]{article}
% Standalone preamble for web-only supplementary notes.
\usepackage[T1]{fontenc}
\usepackage{amsmath,amssymb,amsthm}
\usepackage{hyperref}
\usepackage{cleveref}
\DeclareMathOperator{\Map}{Map}
\DeclareMathOperator{\Fun}{Fun}
\newcommand{\iso}{\xrightarrow{\sim}}
\newtheorem{lemma}{Supplementary Lemma}
\newtheorem{corollary}[lemma]{Supplementary Corollary}
\crefname{lemma}{Supplementary Lemma}{Supplementary Lemmas}
\Crefname{lemma}{Supplementary Lemma}{Supplementary Lemmas}
\crefname{corollary}{Supplementary Corollary}{Supplementary Corollaries}
\Crefname{corollary}{Supplementary Corollary}{Supplementary Corollaries}

\begin{document}
\section*{Pairing and precomposition}

The proof of \href{equivalences.html\#lem:Finite_Compatibility_Of_Products}{pairing and precomposition}
reduces several calculations to the two projections of a product. This note
spells out that reduction and the comparisons it must retain.
All categories and functors in the statements below are absolute.
We write $\beta_1(f,g)\colon\pi_1\circ(f,g)=f$ and
$\beta_2(f,g)\colon\pi_2\circ(f,g)=g$ for the chosen projection
identifications, and
\[
 \kappa(f,g,r)\colon(f,g)\circ r=(f\circ r,g\circ r)
\]
for the precomposition comparison. Vertical composition $\delta\cdot\gamma$
applies $\gamma$ first. The symbol $\star$ denotes whiskering.

\begin{lemma}
\label[lemma]{lem:Supplementary_Pairing_Detection}
Let $h,k\colon X\to C\times D$ and $\sigma,\tau\colon h=k$.
Identifications
\[
 \pi_1\star\sigma=\pi_1\star\tau,
 \qquad \pi_2\star\sigma=\pi_2\star\tau
\]
give an identification $\sigma=\tau$.
\end{lemma}
\begin{proof}[Proof idea]
The product axiom gives an equivalence
\[
 (h=k)\longrightarrow
 (\pi_1\circ h=\pi_1\circ k)\times(\pi_2\circ h=\pi_2\circ k)
\]
whose components are the two postwhiskering functors. Pair the given
identifications, using the precomposition comparison to identify the
images of $\sigma$ and $\tau$ with their paired projections, and reflect
the resulting identification through this equivalence. The conclusion is
an identification between $\sigma$ and $\tau$, not a strict equality of
expressions.
\end{proof}

\begin{lemma}
\label[lemma]{lem:Supplementary_Pairing_Composition}
The chosen pairing of identifications satisfies
\[
 (\mathrm{id}_f,\mathrm{id}_g)=\mathrm{id}_{(f,g)},
\]
and, for composable identifications $\sigma_1,\sigma_2$ and
$\tau_1,\tau_2$,
\[
 (\sigma_2\cdot\sigma_1,\tau_2\cdot\tau_1)
   =(\sigma_2,\tau_2)\cdot(\sigma_1,\tau_1).
\]
\end{lemma}
\begin{proof}[Proof idea]
For $\sigma\colon f=f'$ and $\tau\colon g=g'$, the first projection
of the chosen pairing is
\[
 \pi_1\star(\sigma,\tau)
 =\beta_1(f',g')^{-1}\cdot
       \bigl(\sigma\cdot\beta_1(f,g)\bigr).
\]
Thus the first projection of the identity assertion follows by cancelling
$\beta_1(f,g)$ with its inverse. For composition, substitute the displayed
formula twice, reassociate, and cancel the two occurrences of the
intermediate projection comparison. The result is precisely the formula
for $\sigma_2\cdot\sigma_1$. Postwhiskering preserves vertical
composition, so this computes the projection of the right-hand side as
well. Repeat with the second projection and apply
\Cref{lem:Supplementary_Pairing_Detection}.
\end{proof}

\section*{Changing the inputs}

\begin{lemma}
\label[lemma]{lem:Supplementary_Pairing_Naturality}
Let $f,f'\colon X\to C$, $g,g'\colon X\to D$, and
$r,s\colon R\to X$. For identifications $\sigma\colon f=f'$,
$\tau\colon g=g'$, and $\gamma\colon r=s$, there is an identification
\[
 (\sigma\ast\gamma,\tau\ast\gamma)\cdot\kappa(f,g,r)
 =\kappa(f',g',s)\cdot\bigl((\sigma,\tau)\ast\gamma\bigr),
\]
where $\sigma\ast\gamma=(\sigma\star s)\cdot(f\star\gamma)$
is horizontal composition.
\end{lemma}
\begin{proof}[Proof idea]
First hold $r$ fixed and vary $f,g$. Then hold $f,g$ fixed and vary
$r$. For either square, postcompose with a projection and then with the
target projection comparison. The projection formulas for pairing and
$\kappa$ reduce the desired square to the corresponding whiskering and
associator comparisons. Cancel the target comparison and use
\Cref{lem:Supplementary_Pairing_Detection}.

Paste these two squares. The composition formula in
\Cref{lem:Supplementary_Pairing_Composition} combines their paired sides
into $(\sigma\ast\gamma,\tau\ast\gamma)$, and the other sides combine
into the displayed horizontal composite. This argument retains the
chosen $\kappa$ at both ends.
\end{proof}

\section*{Successive substitution}

Write $\alpha_{r,s,f}\colon(f\circ s)\circ r=f\circ(s\circ r)$
for the associator and $\rho_f\colon f\circ\mathrm{id}=f$ for the
right unitor.

\begin{lemma}
\label[lemma]{lem:Supplementary_Pairing_Substitution}
For $f\colon X\to C$, $g\colon X\to D$, $s\colon R\to X$,
and $r\colon Q\to R$, the two identifications
from $((f,g)\circ s)\circ r$ to
$(f\circ(s\circ r),g\circ(s\circ r))$ satisfy
\[
\begin{aligned}
 &\kappa(f,g,s\circ r)\cdot\alpha_{r,s,(f,g)}\\
 &\quad=(\alpha_{r,s,f},\alpha_{r,s,g})\cdot
   \bigl(\kappa(f\circ s,g\circ s,r)\cdot
                    (\kappa(f,g,s)\star r)\bigr).
\end{aligned}
\]
For the identity substitution one has
\[
 (\rho_f,\rho_g)\cdot\kappa(f,g,\mathrm{id})=\rho_{(f,g)}.
\]
\end{lemma}
\begin{proof}[Proof idea]
Project the two substitution paths. The projection formulas for $\kappa$
leave the same projection comparison, transported through two different
strings of associators. The pentagon identifies those strings. After
cancelling the common comparison, apply
\Cref{lem:Supplementary_Pairing_Detection}.

The identity case uses the compatibility of the right unitor with a
composite:
\[
 \rho_{\pi\circ h}
   =(\pi\star\rho_h)\cdot\alpha_{\mathrm{id},h,\pi}.
\]
This compatibility follows from the pentagon and triangle, followed by
cancellation after prewhiskering by the identity. Insert it into each
projected identity square, cancel its projection comparisons, and again
reflect through the product equivalence.
\end{proof}

These are comparisons for the specified pairing and substitution
operations. They do not replace those operations by strictly associative
or strictly unital ones.
\end{document}
